%=============================================================
% tmfeo3_qafm_e12_crosspeak_modeling.tex
% Diagnosis of the (omega_qAFM, omega_E12) cross-peak search in the
% chi2x-driven Route A 2DCS campaign, and a modeling roadmap.
% Companion to util/crosspeak_routeA.py and the
% build/workflow/v3_grid_chi2x075_wline campaign.
%=============================================================
\documentclass[11pt,a4paper]{article}
\usepackage[margin=2.5cm]{geometry}
\usepackage{amsmath,amssymb,bm,mathtools,booktabs,array,tabularx,multirow}
\usepackage{graphicx}
\usepackage{xcolor}
\usepackage{hyperref}
\usepackage{enumitem}
\usepackage{caption}
\usepackage[T1]{fontenc}

\hypersetup{colorlinks=true,linkcolor=blue,citecolor=blue,urlcolor=blue}

% ---- notation shortcuts ----
\newcommand{\vlam}{\bm{\lambda}}
\newcommand{\vS}{\bm{S}}
\newcommand{\Heff}{\mathcal{H}}
\newcommand{\FeOs}{TmFeO$_3$}
\newcommand{\Gtwo}{\Gamma_2}
\newcommand{\qAFM}{\mathrm{qAFM}}
\newcommand{\qFM}{\mathrm{qFM}}
\newcommand{\EtA}{E_{12}}
\newcommand{\EtB}{E_{13}}
\newcommand{\chitwox}{\chi_{2x}}
\newcommand{\chitwoz}{\chi_{2z}}
\newcommand{\crossP}{(\omega_{\qAFM},\,\omega_{\EtA})}
\newcommand{\meVtoTHz}{0.241799}
\newcommand{\sF}{\sigma_F}
\newcommand{\sG}{\sigma_G}
\newcommand{\sC}{\sigma_C}

\definecolor{passgreen}{HTML}{27AE60}
\definecolor{failred}{HTML}{E74C3C}
\definecolor{warnora}{HTML}{E67E22}
\definecolor{lightblue}{HTML}{EAF4FB}
\definecolor{lightgreen}{HTML}{D5F5E3}
\definecolor{lightred}{HTML}{FADBD8}

\newcommand{\PASS}{\textcolor{passgreen}{\textbf{[OK]}}}
\newcommand{\FAIL}{\textcolor{failred}{\textbf{[FAIL]}}}
\newcommand{\WARN}{\textcolor{warnora}{\textbf{[WARN]}}}

\title{\textbf{The $(\omega_{\qAFM},\,\omega_{\EtA})$ Cross-Peak in \FeOs:\\
Diagnosis of the Route~A Campaign and a Modeling Roadmap}}
\author{ClassicalSpin\_Cpp project notes}
\date{June 2026}

\begin{document}
\maketitle

\begin{abstract}
We set out to confirm or deny a nonlinear 2DCS cross-peak at
$\crossP \approx (0.906,\,0.500)$~THz produced by the magnetoelastic / bilinear
coupling $\chitwox$ between the Fe qAFM magnon and the Tm $E_{1}\!\to\!E_{2}$
crystal-field transition. After rewriting the analysis to use the correct
observable channel ($\Lambda^1_G = 4\,\overline{\lambda^1_{\text{local}}}$), running
the $\chitwox=0$ baseline (identically zero nonlinear signal), and fixing a
units error in the $\tau$-domain analysis, we conclude: \textbf{at the campaign
value $\chitwox = 0.75$~meV there is no detectable cross-peak at the bare
$\EtA=0.500$~THz}. The reason is physical, not numerical: the static
mean field generated by $\chitwox$ is large compared to the bare crystal-field
splitting and \emph{renormalizes the Tm transition} down to
$\approx 0.074$--$0.147$~THz, so the spectral weight that should land at
$\EtA$ has moved. This document records the full chain of evidence and then
lays out three concrete modeling routes to recover a clean, interpretable
cross-peak.
\end{abstract}

\tableofcontents

%=============================================================
\section{Physical setup and notation}
%=============================================================

\FeOs{} (space group $Pbnm$, No.~62) in the $\Gtwo$ phase carries two coupled
magnetic sectors:
\begin{itemize}[nosep]
  \item \textbf{Fe} (Wyckoff $4b$, $S=5/2$, SU(2), \texttt{spin\_dim}=3): hosts
        the magnon branches $\omega_{\qFM}\approx 0.378$~THz and
        $\omega_{\qAFM}\approx 0.906$~THz.
  \item \textbf{Tm} (Wyckoff $4c$, non-Kramers qutrit, SU(3),
        \texttt{spin\_dim}=8): hosts the crystal-field (CEF) transitions
        $\EtA \approx 0.500$~THz ($e_1=2.068$~meV) and
        $\EtB \approx 1.200$~THz ($e_2=4.963$~meV).
\end{itemize}

The mixed state vector has \texttt{state\_dim}=44: indices $0$--$11$ are the
4 Fe sites $\times$ 3 components; indices $12$--$43$ are the 4 Tm sites
$\times$ 8 Gell-Mann components $\lambda^1\dots\lambda^8$.

\paragraph{Units (critical).}
The integrator works in natural units $\hbar=1$ with energies in meV:
\begin{itemize}[nosep]
  \item Time is measured in $\hbar/\text{meV}$; one code time unit $=0.6582$~ps.
  \item A frequency returned by \texttt{numpy.fft.fftfreq(N,\,dt)} with
        \texttt{dt} in code units is therefore in \textbf{meV}.
  \item Convert to THz with $f_{\text{THz}} = f_{\text{meV}}\times\meVtoTHz$.
\end{itemize}
The target cross-peak in code units is
$\omega_{\qAFM}=3.748$~meV, $\omega_{\EtA}=2.068$~meV.

%=============================================================
\section{Route A: the intended mechanism}
%=============================================================

The proposed (``Route A'') mechanism for the cross-peak is a two-step
nonlinear cascade driven by the bilinear Fe--Tm coupling $\chitwox$:
\begin{enumerate}[nosep]
  \item Pump~1 (THz, along the appropriate polarization) launches the Fe
        \textbf{qAFM} magnon. This modulates the Fe order parameter at
        $\omega_{\qAFM}$.
  \item Through $\chitwox$, the oscillating Fe moment drives the Tm
        $\sigma_F$ even channel ($\lambda^2$, sublattice-uniform), which in
        turn shifts the Tm $\lambda^1$ coherence at the CEF frequency
        $\EtA$.
  \item In a 2DCS scan over inter-pulse delay $\tau$, the qAFM phase imprinted
        during step~1 should appear as modulation of the $\EtA$ response,
        producing a peak at $\crossP$.
\end{enumerate}

\paragraph{Observable channel (confirmed).}
The nonzero Tm dipole channel in the $\Gtwo$ ground state is the global
$G$-type $\lambda^1$, which equals four times the site-averaged local-frame
$\lambda^1$:
\begin{equation}
  \Lambda^1_G \;=\; \sum_{k=0}^{3}\lambda^1_{\text{local}}(k)
  \;=\; 4\,\overline{\lambda^1_{\text{local}}}
  \;=\; 4\,\big[\texttt{M\_local\_SU3}\big]_{:,0}.
\end{equation}
Empirically (at $\tau=-100$, $t=t_0$, $\chitwox=0.75$):
$\lambda^1_{\text{local}}(k) = [-9.664,-9.664,-9.699,-9.699]\times10^{-6}$
(uniform $\sigma_F$ sign pattern), giving $\Lambda^1_G = -3.87\times10^{-5}$,
while $\Lambda^1_C \approx 6.9\times10^{-8}\approx 0$.

\begin{center}\fcolorbox{warnora}{lightred}{\parbox{0.95\linewidth}{
\WARN{} \textbf{Naming bug in older notes.} The earlier ``$\Lambda^1_C$''
that was claimed to carry the cross-peak is a mislabel. The nonzero channel is
$\Lambda^1_G = 4\,\overline{\lambda^1_{\text{local}}}$; the true $\Lambda^1_C$
is $\approx 0$ by symmetry. The global-frame reader channel
\texttt{M\_global\_SU3}$_{:,0}\approx 0$ averages to zero and yields only a
noise-level spectrum --- this is why the precomputed
\texttt{M\_NL\_SU3\_FF\_$\lambda$1.npy} files are useless for this search.
}}\end{center}

%=============================================================
\section{What we did and what we found}
%=============================================================

\subsection{Baseline: $\chitwox=0$ gives identically zero nonlinear signal}
The run \texttt{v3\_grid\_chi2x/v3\_grid\_chix\_0\_W\_0} produces
$\mathrm{NL}[\lambda^1]=\mathrm{NL}[\lambda^2]=0$ exactly (rms $=0.000$).
This confirms (a) the nonlinear subtraction
$\mathrm{NL}=M_{01}-M_0-M_1$ is clean, and (b) all
\texttt{v3\_grid\_chi2x} campaigns have $\chitwoz=0$ (no residual even coupling).

\subsection{$\chitwox=0.75$, $W=0$: the 2D spectrum}\label{sec:chi2x075}
Using \texttt{util/crosspeak\_routeA.py} on the 4.2~GB
\texttt{v3\_grid\_chix\_0p75\_W\_0} file, with correctly
calibrated \textbf{linear} THz axes (\S\ref{sec:bugs}) and
Hann window (default after the fix):

\begin{center}
\begin{tabular}{lccc}
\toprule
Position $(\omega_\tau,\omega_t)$ [THz, linear] & Channel & Amplitude & S/N \\
\midrule
$(0,\; f_2=0.147)$ DC-in-$\tau$ stripe & $\lambda^1_{\rm loc}$ & $2.39\times10^{-3}$ & 23.4 \\
$(\qAFM=0.906,\;f_2=0.147)$ dressed cross-peak & $\lambda^1_{\rm loc}$ & $\sim 10^{-2}$(a) & $\sim$4--7(a) \\
$(\qAFM,\;\EtA=0.500)$ bare CEF target & $\lambda^1_{\rm loc}$ & $9.96\times10^{-7}$ & $<0.1$ \\
Background rms & $\lambda^1_{\rm loc}$ & $1.02\times10^{-4}$ & --- \\
\midrule
$(0,\; f_1=0.074)$ global-frame DC & $\lambda^1_{\rm FF}$ (reader) & --- & 137 \\
$(\qAFM,\; f_1=0.074)$ & $\lambda^1_{\rm FF}$ & --- & 0.14 \\
$(\qAFM,\;\EtA=0.500)$ & $\lambda^1_{\rm FF}$ & --- & $<0.01$ \\
\bottomrule
\multicolumn{4}{l}{\small (a) Hann or no-apod window required; Gaussian window
$\gamma=0.03$ destroys this peak (see \S\ref{sec:bugs}).}
\end{tabular}
\end{center}

\noindent\FAIL{} The target $\crossP$ sits \emph{at} the noise floor
(S/N~$=1.1$). The only genuine feature is the Fe magnon diagonal at
$(\qAFM,\qAFM)$ (S/N~$=4.3$). The global maximum is a ridge that is
\emph{flat in $\omega_\tau$} at $\omega_t=\omega_{\qAFM}$ --- the signature of a
signal that is constant in $\tau$ (``DC-in-$\tau$''), whose FFT leaks across all
$\omega_\tau$.

\subsection{The key discovery: $\chitwox$ renormalizes the Tm transition}
Inspecting the nonlinear $\lambda^1$ ($=\lambda^1_{\text{local}}$) response in
the time domain at fixed $\tau$, the dominant oscillation is \emph{not} at
$\EtA=0.500$~THz but at
\begin{equation}
  f_1 \approx 0.074~\text{THz} \;(0.306~\text{meV}),\qquad
  f_2 \approx 0.147~\text{THz} \;(0.608~\text{meV}),\qquad f_2/f_1\approx 1.99 .
\end{equation}
The amplitudes at $f_1,f_2$ exceed those at $\EtA$ and $\qAFM$ by
$250$--$500\times$. The origin is a large static mean field: with
$\chitwox=0.75$~meV and $\langle S\rangle\sim$ several,
\begin{equation}
  \chitwox\,\langle S^2\rangle \;\sim\; 6~\text{meV}
  \;\gg\; e_1 = 2.07~\text{meV} = \hbar\omega_{\EtA},
\end{equation}
so the Tm two-level subspace is strongly dressed and its effective splitting
collapses to $\approx f_1$. Consequently the spectral weight that ``should''
appear at $\EtA$ has physically moved to $f_1$, and the 2D FFT shows
\emph{nothing} at the bare $\EtA$.

\subsection{The equilibrium itself is not stationary}
The unpumped reference $M_0$ already oscillates: $\lambda^1$ of the reference
trajectory rings persistently at $f_1,f_2$ with no decay (early/late amplitude
ratio $\approx 1.07$ over the full $-65.8$ to $+65.8$~ps window), while the
populations ($\lambda^3$, $\lambda^8$) stay constant. This means the Tm qutrit
was \textbf{initialized in the bare CEF eigenbasis, not in the
$\chitwox$-dressed eigenbasis}; the mismatch launches a coherent ringing at the
dressed frequency. This ringing contaminates the nonlinear subtraction and is
the source of the strong $f_1,f_2$ lines.

\subsection{Units error that previously masked the picture}
An earlier $\tau$-FFT used \texttt{stride}=4 (effective
$d\tau_{\text{eff}}=0.2$ code units), giving Nyquist $=2.5$~meV~$=0.605$~THz.
Because $\omega_{\qAFM}=3.748$~meV is \emph{above} this, the qAFM modulation was
\textbf{aliased}, and a ``dominant $0.609$~THz tau-oscillation'' was reported
that was actually $0.608$~meV~$=0.147$~THz~$=f_2$. With \texttt{stride}=1 and
correct meV$\to$THz conversion, the genuine complex amplitude at
$\EtA$ does carry an $\omega_\tau=\omega_{\qAFM}$ Fourier component at
$\sim 8\%$ of its DC value --- real, but spread so broadly in the 2D map that it
never rises above the noise floor.

\subsection{Analysis bugs discovered and fixed}\label{sec:bugs}

Three analysis bugs were found and corrected in \texttt{util/crosspeak\_routeA.py}
after running the campaign. Understanding them is essential for correctly
interpreting all previous S/N numbers.

\paragraph{Bug~1: Angular vs.\ linear THz (critical).}
The \texttt{do\_fft2()} function saved omega axes as
$\omega_{\text{angular}} = f_{\text{code}}\times 2\pi\times\meVtoTHz$ (rad/ps),
but the search constants were defined in \emph{linear} THz
($\nu_{\qAFM}=0.906$, $\nu_{\EtA}=0.500$~THz). The search was therefore
looking at
\[
  \nu_{\text{wrong}} = \frac{\omega_{\text{angular}}}{2\pi}
  \;\approx\; 0.906/(2\pi) \approx 0.144~\text{THz}\approx f_2,
\]
i.e.\ \emph{at the dressed Tm frequency}, not at qAFM! All previous S/N values
at the ``qAFM position'' were actually measuring the dressed-Tm response.
Fix: remove the $2\pi$ factor; save in linear THz.

\paragraph{Bug~2: Flip/index mismatch in $\omega_\tau$.}
The spectrum rows are flipped (\texttt{spec[::-1,:]}, rephasing convention), but
the saved \texttt{omega\_tau} array was the \emph{unflipped} DFT frequency grid.
Row~$i$ of the saved spectrum corresponded to
$\omega_\tau[n_\tau - 1 - i]$, not $\omega_\tau[i]$, causing the positive-quadrant
search to look in the wrong direction.
Fix: save \texttt{otau = otau\_unflipped[::-1]}, so index $i$ maps consistently
to the $i$-th row of the flipped spectrum.

\paragraph{Bug~3: Gaussian apodization destroys the qAFM cross-peak.}
The default Gaussian window has $\sigma_\tau = \tau_{\rm range}/k \approx 24.9$~ps.
The product
\[
  \sigma_\tau \cdot \omega_\qAFM^{\text{code}}
  = 37.76~(\text{code})\times 3.748~(\text{code}^{-1}) \approx 141 \gg 1
\]
places $\omega_\qAFM$ far in the tail of the Gaussian spectral envelope.
The DFT at the $\omega_\qAFM$ bin is dominated by the DC sidelobe,
not the genuine qAFM modulation (integration-by-parts estimate:
$\sim 1/(\omega_\qAFM\,d\tau)\approx 5\%$ of DC versus the $\approx 1\%$ true cross-peak).
Without apodization or with a Hann window (which has zero-mean sidelobes in each
basis direction), the cross-peak at $(\qAFM,\,f_2=0.147)$~THz emerges with
S/N\,$\approx\,4$--$7$.
Fix: default apodization changed from Gaussian($\gamma$=0.03) to Hann;
\texttt{--window \{hann,gaussian,none\}} option added.

\begin{center}\fcolorbox{warnora}{lightred}{\parbox{0.95\linewidth}{
\WARN{} The S/N table in the previous version of this document (with Gaussian
apodization and angular-THz axes) is \textbf{incorrect}. The corrected table
is in \S\ref{sec:chi2x075} above. Old files in
\texttt{build/workflow/routeA\_analysis/omega\_*.npy} store angular THz;
divide by $2\pi$ before use.
}}\end{center}

\subsection{Updated physical picture at $\chitwox=0.75$~meV}
With all bugs fixed:
\begin{itemize}[nosep]
  \item The dominant NL feature is the DC-in-$\tau$ stripe at $\omega_t = f_2 = 0.147$~THz:
        the dressed Tm oscillation is excited only at $\tau=0$ (coincident pumps) and
        thus has no $\tau$ modulation under Gaussian apodization.
  \item With Hann or no-apodization window: a genuine cross-peak is detected at
        $(\omega_\tau, \omega_t) = (0.906, 0.147)$~THz (i.e.\ $\qAFM$ vs.\ dressed~$f_2$)
        with S/N~$\approx 4$--$7$. This confirms the Route~A mechanism is active.
  \item No peak at the bare CEF target $(0.906, 0.500)$~THz in any windowing scheme,
        consistent with strong coupling: $\chitwox\langle S^2\rangle \approx 6$~meV
        $\gg e_1 = 2.07$~meV.
\end{itemize}

%=============================================================
\section{Diagnosis summary}
%=============================================================

\begin{center}\fcolorbox{passgreen}{lightgreen}{\parbox{0.95\linewidth}{
\textbf{Bottom line.}
At $\chitwox=0.75$~meV the Route~A mechanism \emph{is confirmed}:
a cross-peak exists at $(\omega_\qAFM,\,f_2)=(0.906,\,0.147)$~THz at
S/N~$\approx 4$--$7$ (Hann/no-apod window).
However, this is the \emph{dressed} frequency --- the strong mean field
($\chitwox\langle S^2\rangle\approx 6$~meV $\gg e_1=2.07$~meV)
renormalizes the Tm transition from 0.500~THz down to $f_2=0.147$~THz.
A Gaussian apodization window with $\gamma=0.03$ (equivalent to
$\sigma_\tau\cdot\omega_\qAFM\approx 141$) destroys the cross-peak via DC
sidelobe leakage; the fix is to use a Hann window (new default in the script).
The bare-CEF cross-peak at $(0.906,\,0.500)$~THz is absent in all windowing
schemes --- to see it we must work in the perturbative regime
($\chitwox\ll e_1/(\langle S^2\rangle)\approx 0.35$~meV).
}}\end{center}

%=============================================================
\section{Modeling roadmap: how to recover the cross-peak}
%=============================================================

The goal is a 2DCS map in which a peak appears at a \emph{predictable}
$(\omega_\tau,\omega_t)$ that we can attribute to qAFM~$\to E_1E_2$ energy
transfer through $\chitwox$. There are three complementary routes, in order of
increasing effort.

\subsection{Route 1 --- Weak-coupling / perturbative regime (fastest)}
Reduce $\chitwox$ so that the dressing is perturbative,
$\chitwox\langle S^2\rangle \ll e_1$. Then the Tm transition stays at the bare
$\EtA$ and the cross-peak amplitude is, to leading order,
\begin{equation}
  I_{\text{cross}}\;\propto\;\chitwox^{2}\,
  \frac{\partial \langle S_{\qAFM}\rangle}{\partial E_1}\Big|_{\text{pump1}}
  \;\times\;\mu_{12}^{\,2},
\end{equation}
i.e.\ it grows as $\chitwox^2$ while the frequency stays pinned at $\EtA$.
\begin{itemize}[nosep]
  \item \textbf{Action.} Run a small grid
        $\chitwox \in \{0.02,\,0.05,\,0.10,\,0.20\}$~meV at $W=0$.
  \item \textbf{Expectation.} The peak should sit at $(0.906,0.500)$ and its
        intensity should scale as $\chitwox^2$. Confirming the quadratic scaling
        is the cleanest possible proof that the peak is the Route~A cascade and
        not an artifact.
  \item \textbf{Sampling.} Keep $d\tau$ fine enough that
        $\omega_{\qAFM}=3.748$~meV is below Nyquist:
        $d\tau < 1/(2\times3.748)=0.133$ code units. The campaign value
        $d\tau=0.05$ is fine.
\end{itemize}

\subsection{Route 2 --- Dress the reference, then linearize (most physical)}
Keep $\chitwox=0.75$ but \emph{remove the spurious ringing} so the spectrum is
clean at the (renormalized) Tm frequency.
\begin{enumerate}[nosep]
  \item \textbf{Relax to the dressed ground state.} Before launching pumps,
        evolve (or minimize) the coupled Fe+Tm system to its true equilibrium so
        that the reference $M_0(t)$ is stationary. Practically: add a long
        damped pre-roll, or solve the self-consistent mean-field fixed point
        $\partial\Heff/\partial\lambda = 0$ for the Tm qutrit in the static Fe
        field, and initialize the qutrit in \emph{that} eigenbasis.
  \item \textbf{Recompute the dressed transition} $\tilde\omega_{12}$ from the
        Hessian of $\Heff$ at the fixed point (this is the $0.074$--$0.147$~THz
        line). The cross-peak will then appear at
        $(\omega_{\qAFM},\,\tilde\omega_{12})$, which is a legitimate physical
        observable --- just at a renormalized frequency.
  \item \textbf{Verify} by overlaying the linear (1D) Tm spectrum: the single
        pump response must show a peak exactly at $\tilde\omega_{12}$.
\end{enumerate}
This route documents the renormalization rather than fighting it, and is the
correct thing to report if the experimental $\chitwox$ really is $\sim$~meV.

\subsection{Route 3 --- $\chitwoz$-mediated even channel (alternative cascade)}
All current campaigns have $\chitwoz=0$. The physical value is
$\chitwoz\approx 0.6$. The $\chitwoz$ coupling drives the Tm \emph{even}
($\lambda^2$, $G$-type) channel directly and can produce a cross-peak with a
different, possibly cleaner, selection rule.
\begin{itemize}[nosep]
  \item \textbf{Action.} Generate a new campaign with $\chitwoz=0.6$,
        $\chitwox=0$ (isolate the even cascade), then $\chitwox=\chitwoz$
        together.
  \item \textbf{Expectation.} The even channel couples
        $\langle S_{\qAFM}^z\rangle^2$ to the $\lambda^2$ population, giving a
        rectified (DC + $2\omega_{\qAFM}$) drive of the Tm levels; the cross-peak
        structure differs from Route~A and should be compared.
\end{itemize}

\subsection{Analysis-side fixes (apply to all routes)}
\begin{itemize}[nosep]
  \item \textbf{Omega axes: always use linear THz.}
        Multiply \texttt{numpy.fft.fftfreq} output by \texttt{MEV\_TO\_THZ}
        (not $2\pi\times$\texttt{MEV\_TO\_THZ}). Label axes
        $\nu=\omega/(2\pi)$ (THz, cycles/ps).
        \PASS{} Fixed in \texttt{do\_fft2()} (see \S\ref{sec:bugs}).
  \item \textbf{Flip/index consistency.} After the rephasing flip
        \texttt{spec[::-1,:]}, save the \emph{reversed} \texttt{omega\_tau}
        so that row $i$ maps to \texttt{omega\_tau[i]}.
        \PASS{} Fixed: \texttt{otau = otau\_unflipped[::-1]}.
  \item \textbf{Use Hann (or no) apodization in $\tau$.}
        Gaussian($\gamma=0.03$) with
        $\sigma_\tau\cdot\omega_\qAFM\approx 141$ kills the cross-peak.
        Hann window gives S/N~$\approx 4$ at the qAFM cross-peak.
        \PASS{} Fixed: new default \texttt{--window hann} in the script.
  \item \textbf{Remove DC-in-$\tau$ before the 2D FFT.} Subtract the
        $\tau$-mean at each $t$:
        $\widetilde{\mathrm{NL}}(\tau,t)=\mathrm{NL}(\tau,t)-\overline{\mathrm{NL}}(\cdot,t)$.
        This kills the flat ridge that dominates the raw map.
  \item \textbf{Use \texttt{stride}=1 in $\tau$} (or verify $d\tau<0.133$) so
        $\omega_{\qAFM}$ is not aliased.
  \item \textbf{Report S/N against the local rms}, and treat anything below
        $\sim 3\sigma$ as absent.
\end{itemize}

%=============================================================
\section{Concrete next experiment}
%=============================================================

The recommended first step is \textbf{Route 1}, because it is decisive and
cheap:
\begin{enumerate}[nosep]
  \item Run the $\chitwox\in\{0.02,0.05,0.10,0.20\}$ grid, $W=0$, $\chitwoz=0$,
        $d\tau=0.05$, same $t$-grid as the current campaign.
  \item For each, run \texttt{crosspeak\_routeA.py} \emph{with the new
        DC-in-$\tau$ subtraction}.
  \item Plot $I(0.906,0.500)$ vs.\ $\chitwox$ on log--log axes.
        A slope of $2$ confirms the Route~A $\chitwox^2$ cascade and a peak
        pinned at the bare $\EtA$.
\end{enumerate}
If the peak frequency \emph{drifts downward} with increasing $\chitwox$ across
this grid, that directly visualizes the renormalization and motivates Route~2
for the strong-coupling endpoint.

\appendix
%=============================================================
\section{File and data inventory}
%=============================================================
\begin{itemize}[nosep]
  \item Analysis script: \texttt{util/crosspeak\_routeA.py}
        (reads \texttt{M\_local\_SU3}, computes 2D NL FFT).
        Updated: linear THz axes, Hann default window, flip fix.
        Options: \texttt{--window \{hann,gaussian,none\}}.
  \item Visualization: \texttt{util/visualize\_2dcs\_chi2x075.py}
        (loads precomputed npy, produces 3-panel figures with linear THz axes).
  \item $\chitwox=0.75$, $W=0$ data:
        \texttt{build/workflow/v3\_grid\_chi2x075\_wline/v3\_grid\_chix\_0p75\_W\_0/}\\
        \texttt{sample\_0/pump\_probe\_spectroscopy.h5} (4.2~GB).
  \item Precomputed reader npy:
        \texttt{M\_NL\_SU3\_FF\_$\lambda$*.npy} (global-frame, angular THz; divide by $2\pi$).
  \item Baseline $\chitwox=0$:
        \texttt{build/workflow/v3\_grid\_chi2x/v3\_grid\_chix\_0\_W\_0/}.
  \item Stored routeA spectra: \texttt{build/workflow/routeA\_analysis/}\\
        \texttt{omega\_*.npy}: \textbf{angular THz (old)} --- divide by $2\pi$ or re-run script.\\
        \texttt{spec\_lambda*.npy}: spectral amplitudes (valid; omega stored separately).\\
        \texttt{fig1\_2d\_spectra.pdf/.png}, \texttt{fig2\_diagnostics.pdf/.png},
        \texttt{fig3\_full\_range.pdf/.png},
        \texttt{fig4\_no\_apod\_hann.png}: diagnostic figures.
\end{itemize}

\section{Lambda-index key (\texttt{M\_local\_SU3})}
\begin{center}
\begin{tabular}{cll}
\toprule
Index & Generator & Meaning / equilibrium value at $\chitwox=0.75$ \\
\midrule
0 & $\lambda^1$ & $E_{12}$ $x$-coherence; $-5.24\times10^{-3}$ \\
1 & $\lambda^2$ & $E_{12}$ $y$-coherence; $-1.35\times10^{-2}$; DC-dominated NL \\
2 & $\lambda^3$ & diagonal population; $0.415$ \\
3--6 & $\lambda^{4..7}$ & $E_{13}/E_{23}$ off-diagonal; $0$ when $\chitwoz=0$ \\
7 & $\lambda^8$ & total population; $0.910$ (constant) \\
\bottomrule
\end{tabular}
\end{center}

\end{document}
